\section{Introduction}
\label{sec:introduction}

Algorithmic redistricting transfers the locus of potential manipulation
from the map-drawing stage (where humans directly choose boundaries) to
the data preparation stage (where the algorithm's inputs are assembled).
An adversary who cannot influence the algorithm's parameters or code
may still be able to achieve partisan outcomes by manipulating the
population counts, geographic boundaries, or demographic tables that
the algorithm uses as inputs.

This paper examines input data manipulation as a gaming vector (V2 in
the R.0 taxonomy~\citep{r0paper}).
We ask: given full access to the data pipeline between the Census Bureau
and the redistricting algorithm, what is the maximum partisan shift an
adversary can achieve without triggering the SHA-256 audit chain?

The question has both empirical and cryptographic components.
The empirical component asks: how much does a given magnitude of data
manipulation shift partisan outcomes?
The cryptographic component asks: at what manipulation magnitude does the
SHA-256 audit chain provide a computationally infeasible barrier to
undetected manipulation?

\subsection{Data Sources}

The \bisect{} algorithm ingests three classes of data:

\begin{enumerate}
\item \textbf{Census P.L.\ 94-171 Redistricting Data}: Population counts
  for each census tract (GEOID), disaggregated by race and Hispanic origin
  for VRA compliance.
  Released every decade by the Census Bureau within 12 months of census day.
\item \textbf{TIGER/Line Shapefiles}: Geographic boundaries for each
  census tract.
  Released annually by the Census Bureau's Geography Division.
  The redistricting-year vintage (e.g., 2020) is the authoritative source.
\item \textbf{ACS 5-Year Tables}: American Community Survey estimates
  of demographic composition, used for VRA boost weight computation.
  Annually released; \bisect{} uses the vintage ending in the census year
  (e.g., 2016--2020 for the 2020 cycle).
\end{enumerate}

All three are public-domain and available directly from the Census Bureau.
The key security property is that they are also independently verifiable:
any party can download them directly and compare against the algorithm's
cached copies.

\subsection{Road Map}

Section~\ref{sec:attack} describes the three attack surfaces.
Sections~\ref{sec:pl94}--\ref{sec:acs} analyse each attack in detail.
Section~\ref{sec:audit-chain} describes the three-level SHA-256 audit chain.
Section~\ref{sec:tls} discloses the TLS limitation.
Section~\ref{sec:effect} quantifies the combined partisan effect.
Section~\ref{sec:conclusion} concludes.
